\chapter{About Mathilda}
\label{ch:about}

\index{Mathilda!goals}%
Mathilda is a computer algebra system---a program\index{computer algebra system} that does mathematics
\emph{symbolically}\index{symbolic computation}, computing with $x$ and $\pi$ and $\sqrt{2}$ as exact objects
rather than as approximate decimals, while remaining fast enough to crunch a
million floating-point numbers when that is what you need. It is modelled closely
on the core architecture and evaluation semantics of the Mathematica\index{Mathematica}, and it
is written in C, and in its own language.

This chapter is about \emph{why} Mathilda exists, what it is built on, and why it
is written the way it is. The next chapter is practical---how to obtain, build, and
run the system on your own machine---and the mathematics begins in earnest in
Chapter~\ref{ch:intro}; here we set the stage.

\section{The mission}
\label{sec:mission}

The goal of Mathilda is simple to state and hard to achieve: \emph{to make the
most advanced, most efficient computer algebra system freely available to the
world.}

Each word in that sentence is load-bearing.

\emph{Advanced}, because a computer algebra system earns its keep at the frontier
---factoring polynomials over algebraic number fields, integrating in closed form
by the Risch algorithm\index{Risch algorithm}, solving Diophantine equations, computing with Gr\"obner\index{Grobner basis@Gr\"obner basis}
bases~\cite{gcl}. Anyone can add two fractions; the interesting question is how far
into real mathematics a system will carry you.

\emph{Efficient}, because a symbolic answer you cannot get before lunch is of
little use. Mathilda is built to be quick: exact arithmetic on arbitrary-precision\index{arbitrary-precision arithmetic}
integers, a bytecode compiler for numeric inner loops\index{bytecode compiler}, and transparently packed\index{array!packed}
arrays that let element-wise work run at the speed of machine code.

\emph{Freely available}, because mathematics belongs to everyone. The dominant
computer algebra systems are proprietary and expensive; their algorithms are
sealed. A student, a researcher in a country without site licences, or a curious
amateur should be able to compute at the highest level without asking anyone's
permission and without paying a toll.

\section{Free software, and why it matters}
\label{sec:license}

\index{license!GPLv3}%
Mathilda is released under the \textbf{GNU General Public License, version~3}
(GPLv3). You may run it for any purpose, read every line of its source, modify it,
and redistribute your modifications---provided you extend the same freedoms to
whomever you pass it on to. This is not an incidental legal choice; it follows
directly from the mission.

For a computer algebra system in particular, openness is a matter of
\emph{mathematical integrity}\index{mathematical integrity}. When a system tells you that
\[
  \int \frac{dx}{1+x^{2}} = \arctan x,
\]
you can check it by differentiating. But when it returns a page-long closed form
for a hard integral, or announces that a two-hundred-digit number is prime, you
are trusting the algorithm. With proprietary software that trust is an act of
faith; with free software it is, in principle, verifiable. The algorithm is right
there to be read, tested, and---when it is wrong---fixed. Scientific results
computed with a black box are not fully reproducible; results computed with an
open system are.

Openness also buys \emph{longevity}\index{longevity}. Proprietary formats and products are
abandoned; a licence server goes dark and a decade of work becomes unreadable.
Free software cannot be taken away. As long as one copy of the source survives, so
does the ability to build it, study it, and carry it forward. A system meant to
serve mathematics---which measures its lifetimes in centuries---should be built to
last, and only free software can make that promise.

\section{Standing on the shoulders of giants}
\label{sec:libraries}

Mathilda does not reinvent the deep numerical kernels that decades of careful
engineering have already perfected. Instead it builds on a small set of
best-in-class libraries, each the state of the art in its niche. You can see them
listed the moment you start the program, in the version banner:

\mtranscript{01-about/version}

Behind that one line stands an enormous amount of work by other people:

\begin{itemize}
  \item \textbf{GMP}~\cite{gmp}\index{GMP}, the GNU Multiple Precision Arithmetic Library, provides the
        arbitrary-precision integers on which all of Mathilda's exact arithmetic
        rests. When a machine integer would overflow, the value is promoted to a
        GMP ``bignum'' automatically and silently.
  \item \textbf{MPFR}~\cite{mpfr}\index{MPFR} extends that idea to real numbers, giving correctly rounded
        arbitrary-precision floating point---so \B{N}\texttt{[expr, 100]} can be
        trusted to a hundred digits.
  \item \textbf{FLINT}~\cite{flint}\index{FLINT} accelerates polynomial arithmetic, factoring, and the
        rigorous numerics behind many special functions.
  \item \textbf{LAPACK}\index{LAPACK} and \textbf{BLAS}\index{BLAS}
        power dense machine-precision linear algebra---the decompositions behind
        \B{Det}, \B{Inverse}, \B{Eigenvalues}, and
        \B{SingularValueDecomposition}.
  \item \textbf{GMP-ECM}\index{GMP-ECM} supplies the Elliptic Curve Method for the harder cases
        of integer factorization.
  \item \textbf{Raylib}\index{Raylib} renders the interactive graphics produced by \B{Plot} and
        \B{Show}, and \textbf{GNU~Readline}\index{GNU Readline} gives the interactive session its
        line editing and history.
\end{itemize}

The art is not in rewriting these libraries but in \emph{orchestrating} them: in
knowing when an exact computation should hand off to GMP, when a numeric one
should drop into LAPACK, and in making those transitions invisible so that you can
think about the mathematics instead of the plumbing. Much of this book is about
that orchestration.

\section{Why Mathilda is implemented in C}
\label{sec:why-c}

A computer algebra system could be written in almost any language. Mathilda is
written in C99\index{C99}, and the choice is deliberate.

\emph{Performance.} C imposes essentially no overhead between your program and the
machine. The libraries Mathilda depends on---GMP, MPFR, FLINT, LAPACK---are
themselves written in C and Fortran, and calling them from C means no marshalling,
no foreign-function boundary, no garbage-collector pauses in the middle of a tight
loop. When Mathilda packs an array of doubles and maps a function over it, the
inner loop is exactly the loop a C programmer would write by hand.

\emph{Language design and predictability.} C is small. Its semantics fit in one
head; its cost model is transparent---you can look at a line and know roughly what
the machine will do. For a system whose entire job is to be trustworthy about
computation, that transparency is a feature, not a limitation.

\emph{Longevity.} C is perhaps the most stable programming language in wide use.
Code written to the C standard decades ago still compiles and runs today, on
hardware that did not exist when it was written. A project that intends to be
useful for a very long time---see \S\ref{sec:license}---should be written in a
language that will still be understood, and still be compilable, a very long time
from now. C is that language.

\emph{An unfair advantage.} There is also a modern, practical reason. Large
language models are exceptionally strong\index{large language model} at writing and reasoning about C: it is
one of the most abundant languages in their training, its patterns are regular,
and its errors are mechanically checkable. Much of Mathilda is developed with such
models in the loop, and C's combination of ubiquity and simplicity makes that
collaboration unusually productive. A language that is easy for both humans and
machines to reason about is exactly the right substrate for building an ambitious
system quickly and carefully.

\section{The long-term vision of Mathilda}
\label{sec:vision}
\index{vision}%

The mission of \S\ref{sec:mission} sets a deliberately immodest bar, and it is worth
being explicit about how high. Mathilda's long-term goal is not merely to be \emph{a}
capable computer algebra system, but to be the \emph{best} system in the world at
every domain of computational mathematics---to gather, in a single free program, the
capabilities that today are scattered across a dozen specialised and often
proprietary tools, and to match or surpass the finest of each on its own ground.

That standard is measured against the best there is, domain by domain. To make it
concrete, here is some of what ``best in each domain'' is meant to mean:

\begin{itemize}
  \item \textbf{Symbolic integration.} A complete implementation of the
        Risch--Trager--Bronstein algorithm\index{Risch--Trager--Bronstein algorithm}
        for integration in closed form, rivalling the Scratchpad / AXIOM / FriCAS\index{FriCAS}
        lineage that pioneered it---so that when an elementary antiderivative exists
        Mathilda finds it, and when none exists Mathilda proves that too.
  \item \textbf{Algebraic geometry and commutative algebra.} Gr\"obner bases,
        ideals, varieties, and the wider machinery of computational algebra at a level
        that rivals Magma\index{Magma}, Singular\index{Singular}, and Macaulay2\index{Macaulay2}.
  \item \textbf{Number theory.} Factorization, primality proving, elliptic curves,
        and class-group computation to rival PARI/GP\index{PARI/GP}.
  \item \textbf{Finite fields and discrete computation.} Arithmetic and linear
        algebra over finite fields at the scale that coding theory and cryptography
        demand, rivalling Magma.
  \item \textbf{Tensor computation.} Manipulation of tensors and the objects of
        differential geometry that is as expressive and as efficient as Cadabra\index{Cadabra}.
  \item \textbf{Arbitrary-precision numerics.} Correctly rounded and rigorously
        bounded computation---ball arithmetic in the spirit of Arb\index{Arb} and mpmath\index{mpmath}---so a
        numerical answer comes with a guarantee, not merely a value.
  \item \textbf{Numerical computing.} Dense linear algebra, quadrature, and
        differential-equation solving at the level of LAPACK, SciPy\index{SciPy}, and Julia\index{Julia}.
  \item \textbf{Optimization.} Local and global optimization---constrained,
        unconstrained, and mixed-integer---to rival SciPy and NLopt\index{NLopt}.
  \item \textbf{Array numerics.} Element-wise and reduction operations on dense,
        packed numerical arrays that are at least as fast as NumPy\index{NumPy} on a single
        core---and, because Mathilda owns its memory layout and kernels end to end,
        that parallelise cleanly across many cores.
  \item \textbf{Compilation.} A bytecode compiler whose numeric throughput rivals
        Pipi, so that an inner loop\index{Pipi} written in Mathilda runs at the speed of compiled
        code rather than interpreted rewriting.
\end{itemize}

\noindent
---and so on, across special functions, polynomial factorization, graph theory, and
every other branch. In each the bar is the same: not ``adequate'' but ``as good as
the best there is,'' and the whole brought together under one roof, freely available.

This is an ambition, stated honestly, not a finished boast. Mathilda does not yet
match every one of those systems on every problem, and this book is candid about where
it falls short: throughout, we hold Mathilda's answers and its speed up against the
best available tools and say plainly where there is still ground to make up. The
vision is the direction of travel. What follows is an honest account of how far along
that road Mathilda has come---and, just as importantly, of \emph{how} it does what it
does, because the point of this book is not only to compute, but to understand the
computation.
